\section{Reflection: What Makes the System Intelligent?}

The system's intelligence is not in any one algorithm but in the interaction of
all ingredients.

\paragraph{Search, optimisation, and representation together.} \textbf{Search} (A*) supplies
global competence --- completeness and optimality over the map, so the robot never traps
itself in a dead end. \textbf{Optimisation} (DWA/MPPI) supplies local competence ---
continuously choosing dynamically feasible motions that trade off progress, safety and
smoothness against the live obstacle geometry. \textbf{Representation} (the probabilistic
risk field and the inflated planning grid) is what lets the planners cooperate: it turns a binary world into a graded one and guarantees the two layers agree on what is
drivable. Intelligent behaviour comes from this combination; removing any one part
degrades the whole (no search $\Rightarrow$ dead ends; no optimisation $\Rightarrow$
undrivable or longer paths; no shared representation $\Rightarrow$ the layers fight).

\paragraph{Adaptation.} The system is reactive. The local optimiser re-solves
its cost minimisation at every control step, and MPPI's receding-horizon warm start carries
information forward in time, so the robot adjusts its command to whatever obstacles are
nearby rather than blindly replaying the global path.

\paragraph{Limitations and improvements.} The environment is static and fully observed at
plan time; dynamic obstacles and true partial observability are not yet handled. DWA needs
manual weight tuning and remains freeze-prone. Natural improvements are: online replanning
of the global path when the map changes, a sampling-based global planner (RRT*) for
continuous spaces, learning the cost weights instead of hand-tuning, and
adding an explicit recovery behaviour. Consistent with the project's emphasis, the value
here is less the specific algorithms than a clear account of why each is needed and
how they combine into intelligent navigation.
